\section{Krav og use cases}

Kravene omsætter problemformuleringen til systemvalg. De definerer, hvad appen skal kunne for at håndtere planlagt træning, faktisk udførelse, completion, progression, lavt betjeningsbesvær og sikker dataadgang. COM-B og BCT Taxonomy gør selvmonitorering, feedback og handlingsstøtte relevante kravspor \cite{michie2011_behaviour_change_wheel,michie2013_bct_taxonomy_v1}, mens effektivt engagement gør hyppig appbrug mindre interessant end brug, der støtter træningshandlingen \cite{yardley2016_effective_engagement_dbci}.

Kravarbejdet fungerer som filter. De centrale krav er dem, der påvirker progression, friktion, datakvalitet eller sikkerhed. Støttefunktioner som import, admin, betaling og trænerrelationer behandles i det omfang, de påvirker disse spor.

\subsection{Kernekrav}

Bilagets use case-overblik viser de bruger- og systemforløb, som kravene bygger på, jf. Tabel~\ref{tab:tm-use-cases}. Fire spor er mest afgørende for problemformuleringen. Appen skal kunne oprette og strukturere programmer, workouts og øvelser. Brugeren skal kunne starte en workout, logge sæt og gemme completion. Cycle runtime skal kunne forbinde gennemført træning med næste relevante handling. watchOS og Live Activity skal fungere som støtteflader til samme træningsmodel.

Kravsporbarheden kobler disse spor til implementeringen, jf. Tabel~\ref{tab:tm-krav-traceability}. F-08 og F-09 dækker tracker og completion. F-13 dækker cyklusprogression. F-15 dækker Live Activity-sporet. Kravene gør ændringer i træningen håndterbare, fordi brugeren kan træne med fuld trackerlog, afslutte uden fuld logging eller fortsætte efter en afbrudt session uden at miste hele programflowet.

Kravene skelner mellem primære og støttende flows. Primære flows er dem, hvor brugeren planlægger, starter, logger, afslutter og fortsætter træning. Støttende flows er for eksempel import, admin, betaling og trænerfunktioner. De støttende flows kan være vigtige for en færdig app, men kerneflowet er afgørende for fleksibel progression.

\subsection{Afvigelser som legitime flows}

Et nøglekrav er, at ændringer i planen behandles som legitime flows. Tracker-on, tracker-off og completion skal kunne adskilles, jf. Figur~\ref{fig:tm-tracker-completion-summary}. Den detaljerede trackerlog giver høj datagranularitet, mens completion-eventet bevarer historik og næste handling, når brugeren ikke logger alle sæt.

ACSM og nyere resistance-training-kilder gør progression til et træningsfagligt krav, mens autoregulering gør variation i faktisk træning relevant \cite{acsm2009_progression_models,currier2026_acsm_resistance_training,greig2020_autoregulation_resistance_training}. Kravene skal derfor kunne rumme både struktur og variation. Bilagets procesmodel viser trackerflowet som en sammenhængende proces, jf. Tabel~\ref{tab:tm-tracker-flow-bpmn}.

Kravlogikken ændrer fejlforståelsen. En afbrudt tracker, en manglende detaljelog eller en ændret øvelse skal ikke automatisk nulstille progressionen. Systemet skal bevare nok information til at give næste handling, samtidig med at det viser, når datagrundlaget er mindre detaljeret. Det er her forskellen på en logbog og en træningsstøtte opstår. Logbogen spørger, hvad der blev registreret. Træningsstøtten spørger også, hvad brugeren bør gøre nu.

\subsection{Lav friktion og sikker adgang}

Lav friktion betyder lavt betjeningsbesvær i selve træningssituationen. Brugeren skal kunne registrere og afslutte træning mellem øvelser, pauser og fysisk belastning. Apple-dokumentationen for Live Activities og Watch Connectivity understøtter brugen af status- og companionflader til tidsfølsom aktivitet \cite{apple_activitykit_live_data_docs,apple_watchconnectivity_docs}. I TræningsMester er pointen færre nødvendige afbrydelser under træning.

Den lette handling skal samtidig have en sikker dataadgang bag sig. OWASP MASVS og Supabase RLS-dokumentationen understøtter princippet om, at mobilklienten ikke bør være eneste adgangsgrænse \cite{owasp2024_masvs,supabase_rls_docs,supabase_securing_api_docs}. Sikkerhedsdiagrammet viser derfor en model, hvor brugerfladen starter handlinger, mens RLS og Edge Functions håndhæver adgang, jf. Figur~\ref{fig:tm-security-boundary-summary}.

Fleksibilitet uden tab af datakontrol kræver samspil mellem brugerflade, state-lag, repositories og backendregler.

\subsection{Samlet kravvurdering}

Kravgrundlaget er stærkest, hvor det direkte forbinder problemformuleringen med kerneflows. Plan, tracker, completion, cycle runtime og næste handling skaber en tydelig linje fra motivation og progression til konkrete use cases og krav. Det giver et analysegrundlag, fordi kravene definerer, hvad appen skal kunne for at håndtere ændringer i træningen fagligt og teknisk.

Kravgrundlaget er mere sekundært for støttefunktioner som import, admin, betaling og trænerrelationer. De kan være nødvendige i en moden app, men de bør vægtes efter deres betydning for datakvalitet, sikker adgang og brugerens mulighed for at fortsætte træningen.

\tmdelkonklusion{Kravets konsekvens}{Kernekravene viser forskellen på en almindelig træningslog og en fleksibel træningsapp. En logbog gemmer det, der skete. TræningsMester skal også kunne bruge det, der skete, til at foreslå næste meningsfulde handling.}
